\begin{center}
    \begin{large}
        \textbf{Abstract}
    \end{large}
    \begin{justify}
        \begin{normalsize}
            This paper will focus on particular group of thermosets,
             namely, polyester resins.
            These liquid resins generally cured to form the matrix
             that surrounds reinforcing materials within the
             structure of composite materials such as
             glass-reinforced plastic.
            This paper initially elucidates the method by which an 
             entirely bio-based polyester resin was manufactured, 
             from simple organic acids and alcohols.
            Subsequently, a description of the method and associated
             results of Bayesion optimisation routines designed to
             improve the strengths exhibited by these materials will
             be brought to bear.
            The most salient findings include the mere possibility
             that these prepolymers can be manufactured simply and
             cheaply.
            Additionally, it is worth noting their strength being
             comparable after optimisation to their contemporary
             incumbents; the styrene-based unsaturated polyester
             resins.
        \end{normalsize}
    \end{justify}
\end{center}